%! TEX program=xelatex
%%%%%%%%%%%%%%%%%%%%%%%%%%%% Document Setup %%%%%%%%%%%%%%%%%%%%%%%%%%%%



% Don't like 10pt? Try 11pt or 12pt
\documentclass[10pt]{article}
\usepackage{fontspec}
\setmainfont{Lato}[
    UprightFont = *-Light,
    BoldFont = *-Bold,
    ItalicFont = *-LightItalic,
    BoldItalicFont = *-BoldItalic
]
\newfontfamily\latosemibold{Lato}[
    UprightFont = *-Regular
]

\usepackage[linkcolor = black]{hyperref}

% The automated optical recognition software used to digitize resume
% information works best with fonts that do not have serifs. This
% command uses a sans serif font throughout. Uncomment both lines (or at
% least the second) to restore a Roman font (i.e., a font with serifs).
% (NOTE: This requires the times package above)
%\renewcommand{\familydefault}{\sfdefault}

% This is a helpful package that puts math inside length specifications
\usepackage{calc}

% This package helps LaTeX auto-hyphenate hyphenated words if you use
% special hyphens. For example, bio\-/mimicry will properly hyphenate
% ``mimicry'' if necessary.
\usepackage[shortcuts]{extdash}

% Layout: Puts the section titles on left side of page
\reversemarginpar

%
%         PAPER SIZE, PAGE NUMBER, AND DOCUMENT LAYOUT NOTES:
%
% The next \usepackage line changes the layout for CV style section
% headings as marginal notes. It also sets up the paper size as either
% letter or A4. By default, letter was used. If A4 paper is desired,
% comment out the letterpaper lines and uncomment the a4paper lines.
%
% As you can see, the margin widths and section title widths can be
% easily adjusted.
%
% ALSO: Notice that the includefoot option can be commented OUT in order
% to put the PAGE NUMBER *IN* the bottom margin. This will make the
% effective text area larger.
%
% IF YOU WISH TO REMOVE THE ``of LASTPAGE'' next to each page number,
% see the note about the +LP and -LP lines below. Comment out the +LP
% and uncomment the -LP.
%
% IF YOU WISH TO REMOVE PAGE NUMBERS, be sure that the includefoot line
% is uncommented and ALSO uncomment the \pagestyle{empty} a few lines
% below.
%

%% Use these lines for letter-sized paper
% \usepackage[paper=letterpaper,
%             %includefoot, % Uncomment to put page number above margin
%             marginparwidth=1.2in,     % Length of section titles
%             marginparsep=.05in,       % Space between titles and text
%             margin=1in,               % 1 inch margins
%             includemp]{geometry}

%% Use these lines for A4-sized paper
\usepackage[paper=a4paper,
            %includefoot, % Uncomment to put page number above margin
            marginparwidth=30.5mm,    % Length of section titles
            marginparsep=1.5mm,       % Space between titles and text
            margin=18mm,              % 25mm margins
            includemp]{geometry}

%% More layout: Get rid of indenting throughout entire document
\setlength{\parindent}{0in}

% Provides special list environments and macros to create new ones
\usepackage[shortlabels]{enumitem}

% Simpler bibsections for CV sections
% (thanks to natbib for inspiration)
%
% * For lists of references with hanging indents and no numbers:
%
%   \begin{bibsection}
%       \item ...
%   \end{bibsection}
%
% * For numbered lists of references (with hanging indents):
%
%   \begin{bibsection}
%       \item ...
%   \end{bibsection}
%
%   Note that bibenum numbers continuously throughout. To reset the
%   counter, use
%
%   \restartlist{bibenum}
%
%   at the place where you want the numbering to reset.

\makeatletter
\newlength{\bibhang}
\setlength{\bibhang}{1em}
\newlength{\bibsep}
 {\@listi \global\bibsep\itemsep \global\advance\bibsep by\parsep}
\newlist{bibsection}{itemize}{3}
\setlist[bibsection]{label=,leftmargin=\bibhang,%
        itemindent=-\bibhang,
        itemsep=\bibsep+2pt,parsep=\z@,partopsep=0pt,
        topsep=0pt}
\newlist{bibenum}{enumerate}{3}
\setlist[bibenum]{label=[\arabic*],resume,leftmargin={\bibhang+\widthof{[999]}},%
        itemindent=-\bibhang,
        itemsep=\bibsep,parsep=\z@,partopsep=0pt,
        topsep=0pt}
\let\oldendbibenum\endbibenum
\def\endbibenum{\oldendbibenum\vspace{-.6\baselineskip}}
\let\oldendbibsection\endbibsection
\def\endbibsection{\oldendbibsection\vspace{-.6\baselineskip}}
\makeatother

\usepackage{xcolor}

%% Reference the last page in the page number
%
% NOTE: comment the +LP line and uncomment the -LP line to have page
%       numbers without the ``of ##'' last page reference)
%
% NOTE: uncomment the \pagestyle{empty} line to get rid of all page
%       numbers (make sure includefoot is commented out above)
%
\usepackage{fancyhdr,lastpage}
\pagestyle{fancy}
%\pagestyle{empty}      % Uncomment this to get rid of page numbers
\fancyhf{}\renewcommand{\headrulewidth}{0pt}
\fancyfootoffset{\marginparsep+\marginparwidth}
\newlength{\footpageshift}
\setlength{\footpageshift}
          {0.5\textwidth+0.5\marginparsep+0.5\marginparwidth-2in}
\lfoot{Last updated: \monthyeardate\today}
\rfoot{\arabic{page}}

\usepackage{datetime}

\newdateformat{monthyeardate}{%
  \monthname[\THEMONTH], \THEYEAR}


% Finally, give us PDF bookmarks
\usepackage{color,hyperref}
\hypersetup{colorlinks,breaklinks,
            linkcolor=black,urlcolor=black,
            anchorcolor=black,citecolor=black}

%%%%%%%%%%%%%%%%%%%%%%%% End Document Setup %%%%%%%%%%%%%%%%%%%%%%%%%%%%


%%%%%%%%%%%%%%%%%%%%%%%%%%% Helper Commands %%%%%%%%%%%%%%%%%%%%%%%%%%%%

%%% HEADING AT TOP OF CURRICULUM VITAE

% The title (name) with a horizontal rule under it
% (optional argument typesets an object right-justified across from name
%  as well)
%
% Usage: \makeheading{name}
%        OR
%        \makeheading[right_object]{name}
%
% Place at top of document. It should be the first thing.
% If ``right_object'' is provided in the square-braced optional
% argument, it will be right justified on the same line as ``name'' at
% the top of the CV. For example:
%
%       \makeheading[\emph{Curriculum vitae}]{Your Name}
%
% will put an emphasized ``Curriculum vitae'' at the top of the document
% as a title. Likewise, a picture could be included:
%
%   \makeheading[{\includegraphics[height=1.5in]{my_picture}}]{Your Name}
%
% the picture will be flush right across from the name. For this example
% to work, make sure the extra set of curly braces is included. Also
% makes ure that \usepackage{graphicx} is somewhere in the preamble.
\newcommand{\makeheading}[2][]%
        {\hspace*{-\marginparsep minus \marginparwidth}%
         \begin{minipage}[t]{\textwidth+\marginparwidth+\marginparsep}%
             {\LARGE \bfseries #2 \hfill #1}%
         \end{minipage}}

%%% SECTION HEADINGS

% The section headings. Flush left in small caps down pseudo-margin.
%
% Usage: \section{section name}
\renewcommand{\section}[1]{\pagebreak[3]%
    \vspace{1.6\baselineskip}%
    \phantomsection\addcontentsline{toc}{section}{#1}%
    \noindent\llap{\latosemibold\smash{\parbox[t]{\marginparwidth}{\hyphenpenalty=10000\raggedright #1}}}%
    \vspace{-\baselineskip}\par}

%%% LISTS

% This macro alters a list by removing some of the space that follows the list
% (is used by lists below)
\newcommand*\fixendlist[1]{%
    \expandafter\let\csname preFixEndListend#1\expandafter\endcsname\csname end#1\endcsname
    \expandafter\def\csname end#1\endcsname{\csname preFixEndListend#1\endcsname\vspace{-0.6\baselineskip}}}

% These macros help ensure that items in outer-type lists do not get
% separated from the next line by a page break
% (they are used by lists below)
\let\originalItem\item
\newcommand*\fixouterlist[1]{%
    \expandafter\let\csname preFixOuterList#1\expandafter\endcsname\csname #1\endcsname
    \expandafter\def\csname #1\endcsname{\let\oldItem\item\def\item{\pagebreak[2]\oldItem}\csname preFixOuterList#1\endcsname}
    \expandafter\let\csname preFixOuterListend#1\expandafter\endcsname\csname end#1\endcsname
    \expandafter\def\csname end#1\endcsname{\let\item\oldItem\csname preFixOuterListend#1\endcsname}}
\newcommand*\fixinnerlist[1]{%
    \expandafter\let\csname preFixInnerList#1\expandafter\endcsname\csname #1\endcsname
    \expandafter\def\csname #1\endcsname{\let\oldItem\item\let\item\originalItem\csname preFixInnerList#1\endcsname}
    \expandafter\let\csname preFixInnerListend#1\expandafter\endcsname\csname end#1\endcsname
    \expandafter\def\csname end#1\endcsname{\csname preFixInnerListend#1\endcsname\let\item\oldItem}}

% An itemize-style list with lots of space between items
%
% Usage:
%   \begin{outerlist}
%       \item ...    % (or \item[] for no bullet)
%   \end{outerlist}
\newlist{outerlist}{itemize}{3}
    \setlist[outerlist]{label=\enskip\textbullet,leftmargin=*}
    \fixendlist{outerlist}
    \fixouterlist{outerlist}

% An environment IDENTICAL to outerlist that has better pre-list spacing
% when used as the first thing in a \section
%
% Usage:
%   \begin{lonelist}
%       \item ...    % (or \item[] for no bullet)
%   \end{lonelist}
\newlist{lonelist}{itemize}{3}
    \setlist[lonelist]{label=\enskip\textbullet,leftmargin=*,partopsep=0pt,topsep=0pt}
    \fixendlist{lonelist}
    \fixouterlist{lonelist}

% An itemize-style list with little space between items
%
% Usage:
%   \begin{innerlist}
%       \item ...    % (or \item[] for no bullet)
%   \end{innerlist}
\newlist{innerlist}{itemize}{3}
    \setlist[innerlist]{label=\enskip\textbullet,leftmargin=*,parsep=0pt,itemsep=0pt,topsep=0pt,partopsep=0pt}
    \fixinnerlist{innerlist}

% An environment IDENTICAL to innerlist that has better pre-list spacing
% when used as the first thing in a \section
%
% Usage:
%   \begin{loneinnerlist}
%       \item ...    % (or \item[] for no bullet)
%   \end{loneinnerlist}
\newlist{loneinnerlist}{itemize}{3}
    \setlist[loneinnerlist]{label=\enskip\textbullet,leftmargin=*,parsep=0pt,itemsep=0pt,topsep=0pt,partopsep=0pt}
    \fixendlist{loneinnerlist}
    \fixinnerlist{loneinnerlist}

%%% EXTRA SPACE

% To add some paragraph space between lines.
% This also tells LaTeX to preferably break a page on one of these gaps
% if there is a needed pagebreak nearby.
\newcommand{\blankline}{\quad\pagebreak[3]}
\newcommand{\halfblankline}{\quad\vspace{-0.5\baselineskip}\pagebreak[3]}

%%% FORMATTING MACROS

% Provides a linked \doi{#1} that links doi:#1 to http://dx.doi.org/#1
%\usepackage{doi}
% To change the text before the DOI, adjust this command
%\renewcommand\doitext{doi:}

% Provides a linked \url{#1} that doesn't require escape characters
\usepackage{url}

% You can adjust the style \url{} uses here:
% (options are: same, rm, sf, tt; defaults to tt)
\urlstyle{same}

% For \email{ADDRESS}, links ADDRESS to the url mailto:ADDRESS
% (uncomment to typeset the e\-/mail address in typewriter font;
%  otherwise, will be typeset in the \urlstyle above)
%\DeclareUrlCommand\emaillink{\urlstyle{tt}}
\providecommand*\emaillink[1]{\nolinkurl{#1}}
\providecommand*\email[1]{\href{mailto:#1}{\emaillink{#1}}}

\providecommand\BibTeX{{B\kern-.05em{\sc i\kern-.025em b}\kern-.08em \TeX}}
\providecommand\Matlab{\textsc{Matlab}}

% Custom hyphenation rules for words that LaTeX has trouble with
\hyphenation{bio-mim-ic-ry bio-in-spi-ra-tion re-us-a-ble pro-vid-er Media-Wiki}
\usepackage{stmaryrd}

%%%%%%%%%%%%%%%%%%%%%%%% End Helper Commands %%%%%%%%%%%%%%%%%%%%%%%%%%%

%%%%%%%%%%%%%%%%%%%%%%%%% Begin CV Document %%%%%%%%%%%%%%%%%%%%%%%%%%%%
% !TeX spellcheck = en_US 
\begin{document}
\makeheading{Manu Navjeevan}

\blankline

\section{Contact Information}

\begin{tabular}{@{}l @{\hskip0.9in}l}
Texas A\&M Department of Economics      &  Email: \href{mailto:mnavjeevan@tamu.edu}     {mnavjeevan@tamu.edu} \\
2935 Research Pw, Office 244 &   Web: \href{https://navjeevan.dev}{https:/\!\!/\! navjeevan.dev}\\
College Station, TX 77843 \\
\end{tabular}

% NOTE: Mind where the & separators and \\ breaks are in the following
%       table. Table is one row made up of three parboxes. The left
%       parbox has address info, the middle parbox has a vertical bar,
%       and the right parbox has phone and electronic contact
%       information.
%
% MACROS: \rcollength is the width of the right column of the table
%             (  it to your liking; default is 1.85in).
%         \spacewidth is width of area between left and right boxes.
%
%\newlength{\rcollength}\setlength{\rcollength}{1.85in}%
%\newlength{\spacewidth}\setlength{\spacewidth}{10pt}
%
%\begin{tabular}[t]{@{}p{\textwidth-\rcollength-\spacewidth}@{}p{\spacewidth}@{}p{\rcollength}}%

% Address box
%\parbox{\textwidth-\rcollength-\spacewidth}{%
%Universidad Nacional de La Plata, FCE\\
%Calle 6, 777 La Plata, Argentina \\
% \\
% \\
%}

%&
% Uncomment to add a vertical bar in middle of contact information
%{\vrule width 0.5pt}
%\parbox[m][5\baselineskip]{\spacewidth}{} &

% Non-snail-mail contact information
%\parbox{\rcollength}{%
%\textit{Phone:} +54-221-6071133 \\
%\textit{E-mail:} \email{serranojqn@gmail.com} \\
% \\
% \\
%}

%\end{tabular}

%%
%% In modern CV's, it seems like ``Objective'' is frowned upon. Instead,
%% incorporate it into a well-constructed cover letter. The ``More
%% information'' can go at the end of the CV, but it should not distract
%% from the section giving references available to contact.
%%
%
% \section{Objective}
%
% Placement in an academic position (i.e., faculty, postdoctoral, or
% research scientist) that allows for advanced research in distributed
% complex adaptive systems (i.e., modeling, analysis, design, and
% verification) with a particular focus on the control of engineered
% agents (e.g., for communications, control, software, electronics, and
% sustainability) and the analysis of biological phenomena (e.g.,
% self-organization, ecological rationality)
% \begin{innerlist}
% \item More information and auxiliary documents can be found at\\\url{http://www.tedpavlic.com/facjobsearch/}
% \end{innerlist}

\section{Academic Appointments}

Texas A\&M University
\begin{outerlist}
    \item[] Assistant Professor of Economics, 2024--
\end{outerlist}

\section{Education}


University of California, Los Angeles
% , Los Angeles, California
\begin{outerlist}
\item[] Ph.D.,
    {Economics}

    % \begin{innerlist}
         % \item[]
    Advisers: {Denis Chetverikov (Primary), Andres Santos, Zhipeng Liao}
     % \end{innerlist}
            
\end{outerlist}


\halfblankline

Carnegie Mellon University
% , Pittsburgh, Pennsylvania
\begin{outerlist}

\item[] B.S.,
        {Economics and Mathematical Sciences}
        % \begin{innerlist}
         
        Graduated with University and College Honors
        % \end{innerlist}
       % \begin{innerlist}
        %\item[] Thesis: \emph{Economic cycle and deceleration of female labor force participation in Latin America}
        %\end{innerlist}


\end{outerlist}


\section{Research \texorpdfstring{\newline}{} Interests}

Econometrics, Causal Inference, Machine Learning
%
% \textbf{High-Dimensional Econometrics:} Orthogonal Learning, Nonparametric Estimation, Post-Selection Inference, High-Dimensional Weak Identification, Effecient Inference
%
% \halfblankline
%
% \textbf{Causal Inference:} Instrumental Variables, Identification with Multiple Instruments, Generalized Monotonicity Conditions
%
\halfblankline

\halfblankline

\halfblankline

\section{Publications}

\begin{bibsection}
    \item \href{https://arxiv.org/abs/2301.06283}{Doubly-Robust Inference for Conditional Average Treatment Effects with High-Dimensional Controls} (with Adam Baybutt) --- \textit{Journal of Econometrics}, 2026
\end{bibsection}

\section{Working \texorpdfstring{\newline}{} Papers}

\begin{bibsection}

    \item Inference under First-Order Degeneracy (with Xinyue Bei)
    \item An Identification and Dimensionality Robust Test for Instrumental Variables Models --- \textit{Revise and Resubmit, Journal of Econometrics}

    \item Identification in Instrumental Variables Models: The Central Role of Abadie's Kappa (with Rodrigo Pinto and Andres Santos) --- \textit{Conditionally Accepted, Econometrica}

    \item \href{https://navjeevan.dev/files/Papers/Ord_Unord_Paper.pdf}{Ordered, Unordered, and Minimal Monotonicity} (with Rodrigo Pinto)

\end{bibsection}

\section{In Progress}
%
\begin{bibsection}
\item On Placebo Testing (with Jorge Garcia and Han Xu)
\item Nonparametric Estimation Under Discrete Covariates with Large Support (with Denis Chetverikov and Andrei Voronin)
\item Compliers, Defiers and the Evaluation of Information Interventions (with Marco Castillo and Mofffii Ige)

\end{bibsection}
\halfblankline

\section{Awards and Fellowships}

Dissertation Year Fellowship, UCLA Graduate Division

\halfblankline

Honors Pass, Econometrics Qualifying Exam, Department of Economics, UCLA

\halfblankline

Honors Pass, Microeconomics Qualifying Exam, Department of Economics, UCLA

% \halfblankline
%
% \textbf{Best Undergraduate Honors Thesis}, Department of Economics, Carnegie Mellon University, 2018






% \section{Submitted Journal Publications}
%
% % Add a little space to nudge next ``Ref'd Journal Publications'' marginpar
% % down to make room for tall ``Submitted Journal Publications''
% % marginpar. If there are enough submitted journal publications, this
% % space will not be needed (and should be removed).
% \vspace{0.1in}




\section{Teaching Experience}
 
Texas A\&M University

\begin{outerlist}
    \item[] Econ 471 (Data Science for Future Decision Makers), 2024--2026
    \item[] Econ 676 (Econometrics II), 2025--2026
\end{outerlist}

\halfblankline

University of California, Los Angeles

\begin{outerlist}
\item[] Econ 103 (Introduction to Econometrics), Summer 2021--2023 
    % \begin{innerlist}
    % \item {My materials from Summer 2021 can be found on \href{https://github.com/mnavjeev/Econ103-Summer-2021}{\underline{GitHub}}}.
    % \end{innerlist}
\end{outerlist}

\section{Refereeing}
American Economic Review, Econometric Reviews,  Econometric Theory, Journal of Business and Economic Statistics, Journal of Econometrics, Journal of Political Economy, Review of Economics and Statistics, Quantitative Economics 

%
% \section{References}
%
% \begin{bibsection}
% \item Prof.\ Denis Chetverikov: \href{mailto:chetverikov@econ.ucla.edu}{chetverikov@econ.ucla.edu} 
% \item Prof.\ Andres Santos: \href{mailto:me@andres@econ.ucla.edu}{andres@econ.ucla.edu} 
% \item Prof.\ Zhipeng Liao: \href{mailto:zihpeng.liao@econ.ucla.edu}{zhipeng.liao@econ.ucla.edu} 
% \end{bibsection}
%


% The ``More Info'' section may not be necessary; make sure it's short
% so it doesn't prevent people from seeing references available to
% contact.
% \section{More Information}

% More information and auxiliary documents can be found at\\%
% \url{http://www.tedpavlic.com/facjobsearch/}.

\end{document}

%%%%%%%%%%%%%%%%%%%%%%%%%% End CV Document %%%%%%%%%%%%%%%%%%%%%%%%%%%%%

%----------------------------------------------------------------------%
% The following is copyright and licensing information for
% redistribution of this LaTeX source code; it also includes a liability
% statement. If this source code is not being redistributed to others,
% it may be omitted. It has no effect on the function of the above code.
%----------------------------------------------------------------------%
% Copyright (c) 2007, 2008, 2009, 2010, 2011 by Theodore P. Pavlic
%
% Unless otherwise expressly stated, this work is licensed under the
% Creative Commons Attribution-Noncommercial 3.0 United States License. To
% view a copy of this license, visit
% http://creativecommons.org/licenses/by-nc/3.0/us/ or send a letter to
% Creative Commons, 171 Second Street, Suite 300, San Francisco,
% California, 94105, USA.
%
% THE SOFTWARE IS PROVIDED "AS IS", WITHOUT WARRANTY OF ANY KIND, EXPRESS
% OR IMPLIED, INCLUDING BUT NOT LIMITED TO THE WARRANTIES OF
% MERCHANTABILITY, FITNESS FOR A PARTICULAR PURPOSE AND NONINFRINGEMENT.
% IN NO EVENT SHALL THE AUTHORS OR COPYRIGHT HOLDERS BE LIABLE FOR ANY
% CLAIM, DAMAGES OR OTHER LIABILITY, WHETHER IN AN ACTION OF CONTRACT,
% TORT OR OTHERWISE, ARISING FROM, OUT OF OR IN CONNECTION WITH THE
% SOFTWARE OR THE USE OR OTHER DEALINGS IN THE SOFTWARE.
%----------------------------------------------------------------------%
